\begin{theorem}[The embedded period curve]\hypertarget{tag-01RE}{}\label{fam:src:F.3.1}\label{fam:st:F.3.1}
\StatementRecord{fam:src:F.3.1}{proved}{01RE}{\SeriesRef{fam:eq:F-j}, \SeriesRef{fam:src:F.2.2}}

For every admissible $c$, the marked period map on the covering of
$B^\circ$ is a holomorphic immersion into $\mathcal O_+$.  The induced
map
\[
 p_c:B^\circ\longrightarrow\Gamma\backslash\mathcal O_+
\]
is injective and a holomorphic immersion.  It extends over $t=0,1$ to a
proper holomorphic embedding
\[
 \bar p_c:\mathbb C=B\setminus\{\infty\}
 \lhook\joinrel\longrightarrow\Gamma\backslash\mathcal O_+.
\]
The closure of $p_c(B^\circ)$ in this quotient is exactly
$\bar p_c(\mathbb C)$.
\end{theorem}



\begin{proof}
\begin{enumerate}[label=\arabic*.,leftmargin=*]
\item For $B=\bigl(\begin{smallmatrix}a&b\\c&d\end{smallmatrix}\bigr)$, changing
the marking by $g(B,x,t)$ replaces $\Pi$ by $\Pi g(B,x,t)^{-1}$ and then
renormalizes its right block.  On the quotient by the complex line of
$e$, the two periods are the row $(\tau,1)$.  Right multiplication by
$B^{-1}$ sends it to $(d\tau-c,a-b\tau)$.  Hence
\begin{equation}\hypertarget{tag-01OX}{}\label{fam:eq:F-tauaction}
 \tau\longmapsto\frac{d\tau-c}{a-b\tau}.
\end{equation}
The denominator never vanishes in the upper half plane.  The function
$j(\tau)$ is therefore a holomorphic invariant of this action and
descends to a holomorphic function $j_\Gamma$ on the proper quotient.
Here $j$ is the elliptic invariant, normalized by
$j(E)=1728\,4A^3/(4A^3+27B^2)$ for $y^2=x^3+Ax+B$; changing a lattice
basis preserves the isomorphism of elliptic curves and hence this
function.

\item We derive the required value in the present family.  The elliptic curve
whose period is $\tau$ has the Engel equation of
Lemma~\ref{I:weierstrass-comparison},
\[
 y^2=x^3-3t^3(t-1)x+2t^4(t-1)^2.
\]
Its discriminant is $1728t^8(t-1)^3$ and direct cancellation gives
\begin{equation}\hypertarget{tag-01OU}{}\label{fam:eq:F-j}
 j(\tau)=1728t\qquad(t\in B^\circ).
\end{equation}
In a local base coordinate $t$, the
derivative of the composite $j\circ\tau$ equals $1728$.  Thus the
derivative of the marked period map is nonzero.  The same identity after
passage to the quotient shows both injectivity and injectivity of the
tangent map of $p_c$: $j_\Gamma/1728$ is a holomorphic left inverse.

\item At the two elliptic points, the period functions are holomorphic in the
cyclic covering coordinates $s_1,s_2$, with $t=s_1^3$ and
$t-1=s_2^4$ in the root coordinates of Definition~\ref{def:common-conventions}.  Their
values are nondegenerate because $D_c<0$ there.  The maps from these
discs to $\mathcal O_+$ are equivariant for $A_1,A_2$.  Composing with
the quotient makes them invariant, so they descend to holomorphic maps
from the base discs.  This is the finite-quotient comparison map: a
holomorphic invariant germ on a disc is a convergent series in $s_j^{m_j}$,
which proves descent on local function rings.  The extensions agree
with $p_c$ off their centers and therefore with each other.  Identity
\eqref{fam:eq:F-j} extends by holomorphy, so the resulting $\bar p_c$ still
has holomorphic left inverse $j_\Gamma/1728$.

\item The quotient is Hausdorff by \ref{fam:src:F.2.2}.  For a compact subset $K$ of it,
$\bar p_c^{-1}(K)$ is closed in $\mathbb C$ and lies in the compact set
$j_\Gamma(K)/1728$.  Thus $\bar p_c$ is proper.  A holomorphic section of
a holomorphic map is locally a closed immersion: in local embeddings
into polydiscs, its coordinate along the base is $t$ and the remaining
coordinates are its holomorphic component functions, whose graph
defines the image.  Applying this to the left inverse shows that
$\bar p_c$ is a holomorphic embedding; properness makes its image
closed.  Since $\mathbb C\setminus\{0,1\}$ is dense in $\mathbb C$, its
image has the asserted closure.
\end{enumerate}
\end{proof}